\chapter*{Free Will and Counting to 3}
\addcontentsline{toc}{chapter}{Free Will and Counting to 3}

The machine begins with a small question: can you tell the difference?

Not can you explain the universe. Not can you derive a constant that
frightened better physicists than either of us. Not can you make a
theory large enough to stand in for reality. The first question is
smaller and more severe. Can a distinction be made, repeated, and
checked?

That is the beginning of this book. A computation becomes an experiment
when it can be repeated by someone who does not trust you. The repeat is
the mercy. It takes the burden off persuasion. The witness does not have
to believe the author, the diagrams, or the mood of discovery around the
instrument. The witness runs the machine again and asks whether the same
reading comes back.

This volume is about how to read that return.

The instrument underneath the story has a deliberately narrow public
face. A claimed calibration enters. The machine compares it against a
repeatable reading and returns whether the claim survived the floor,
together with the residual that remains. That is all a calibration
verifier owes us. It should not ask for reverence. It should not ask for
philosophical permission. It should expose a bounded answer and let the
reader decide what to do next.

The reader's freedom is not removed by the machine. It is sharpened.
Free will, in this book, is the right to choose the next distinction to
ask for. The machine's answer is bounded by what it can count.

\section*{The First Count}

The first count is a reading.

Something has changed. Something is not the same as the blank. This is
the smallest act the instrument can recognize. In ordinary language we
might call it a bit, a mark, a measurement, a signal, or a fact. Those
words come from different histories. Here, before any of that, the first
count means only this:

\[
\hbox{one distinction}\quad\longrightarrow\quad\hbox{a reading}.
\]

A reading is not yet an orientation. It has no left or right. It does
not know whether it is matter or antimatter, inward or outward, before
or after. It is the machine saying: there is enough difference here to
refuse silence.

That refusal is already a kind of freedom. Not human freedom, not moral
freedom, not the old quarrel over whether a mind can begin a causal
chain. It is the smaller freedom of a grammar that has not collapsed
every case into the same case. A machine that cannot distinguish
anything cannot measure anything. A reader who cannot ask for a
distinction cannot verify anything.

So the first count is the first courtesy the instrument pays the reader.
It says: here is something you may check.

\section*{The Second Count}

The second count gives orientation.

Once there are two distinguishable readings, the machine can compare
them. It can say not merely that there is a difference, but that the
difference has a direction. Later chapters read that direction as sign,
charge, spin, residue, phase, or strain. In this preface the claim stays
smaller:

\[
\hbox{two distinctions}\quad\longrightarrow\quad\hbox{an orientation}.
\]

Orientation is the beginning of experiment. A single mark may be an
accident. Two marks can disagree. They can define a before and an after.
They can be swapped, reflected, paired, or cancelled. They can produce a
residual.

The important word is not direction but check. Direction gives the
reader something to test. If a sign flips when the frame flips, that
flip is not poetry. It is a handle. If a residual changes when a channel
is rotated, the machine has exposed a place where the claim can fail.

A failed reading is useful when it lands in a named channel. A vague
failure is fog; a named failure is information. The grammar therefore
needs outcomes that can be read as signals, not moods. Full closure,
structural closure, outgrowth beyond the floor, a leaking ledger, a
broken monotonic refinement, and a sacred-boundary violation are
different events. They must not be flattened into the same embarrassed
word.

\section*{The Third Count}

The third count gives a frame.

With three independent counted directions, the machine can stop leaning
on an imagined continuum. It has enough room to say here, there, and
around. It has a local room in which a reading can be placed. That is
the first complete gauge:

\[
\hbox{three distinctions}\quad\longrightarrow\quad\hbox{a frame}.
\]

This is why counting to three matters. It is not nostalgia for
arithmetic. It is the smallest complete spatial discipline the
instrument can use without borrowing infinity. In the new plan for the
reading, the three directions are written in polar form:

\[
r,\qquad \theta,\qquad \phi .
\]

Those are not used as global coordinates for an ordinary Lorentz
transform. That would be the wrong object. The plan is local. The polar
chart is read into a local frame, a tetrad:

\[
(r,\theta,\phi)
   \longrightarrow (x_r,x_\theta,x_\phi)
   \longrightarrow (t,x_r,x_\theta,x_\phi).
\]

Only there does the signed norm get checked. The frame is local because
the instrument is finite. It does not pretend to own the manifold. It
owns the room it can count.

The strange part, and the part that makes the construction useful for
this book, is the glue. Each spatial leg is also read as a radiation
channel:

\[
\begin{array}{rcl}
r      &\longrightarrow& \hbox{alpha-particle channel},\\
\theta &\longrightarrow& \hbox{beta-particle channel},\\
\phi   &\longrightarrow& \hbox{gamma-particle channel}.
\end{array}
\]

The names are dangerous if they are handled casually. The
alpha-particle channel is not the fine-structure constant. It is a
radial, mass-bearing radiation gauge. The beta-particle channel is the
charged tangent deflection. The gamma-particle channel is the azimuthal
phase or photon-like holonomy. The book will keep those names separate
because a name collision here would fake understanding.

The point of the glue is not to smuggle physics into a coordinate
system. The point is to give the three counted dimensions something to
measure. The local frame gives the spatial count. The three radiation
channels give the residue ledger. If the ledger does not balance, the
reading is not valid.

\section*{The Floor and the Needle}

Every instrument needs a floor.

Below the floor, a distinction may be real in some larger sense, but
this machine is not licensed to claim it. The floor is not
embarrassment. It is the honesty condition. Without a floor, every tail
becomes a prophecy. With a floor, the instrument can say: this much is
above the floor; this much is below the floor; this much is the
residual.

The blank calibration is the clean identity:

\[
\hbox{TRUE}=\hbox{TRUE}.
\]

That is where the instrument learns how much effort it costs to
recognize the cleanest identity. The floor is read before the
interesting claim. If the floor is moved after the claim, the experiment
is no longer blind.

The second blank is

\[
.999 = 1 .
\]

This book does not read that as the infinite decimal identity usually
written with an ellipsis. The ellipsis is not the point. The point is
finite and counted. The current needle counts three dimensions. It
certifies a box:

\[
137.035999 \leq \alpha^{-1} \leq 137.036 .
\]

The machine does not own the limit. It owns this box.

This is the sentence the rest of the volume must protect. It is stronger
than triumph and less brittle. It does not say the machine has captured
every digit. It says the machine has counted enough dimensions to
certify the current interval, and that any sharper claim must pay for
another count.

That payment is not rhetorical. It is not a prettier paragraph. It is
another sensor class, another residue, another run through the verifier.
More precision requires counting more dimensions.

\section*{The Reader's Part}

The reader is not passive in this story.

The reader chooses what claim to submit to the verifier. The reader
chooses whether to trust the setup enough to run it. The reader chooses
whether a structural result is enough for the next question or whether
another dimension must be counted.

This is why free will belongs in the preface. Not as mysticism, and not
as an escape hatch from determinism. The machine is deterministic. That
is what makes it repeatable. The freedom is upstream of the run, in the
choice of which distinction to ask for, and downstream of the run, in
the judgment of what the named outcome licenses.

A computation that always returns the same result can still enlarge
human freedom, because it removes a kind of argument. It lets us stop
asking who sounded more convincing and start asking what survived the
floor.

\section*{What This Book Promises}

This book promises to be a guide to the reading, not a sermon about the
answer.

It will explain the floor, the needle, the polar frame, the three
radiation channels, the search, and the verifier. It will cite the
standard mathematics and physics generously. It will let the text go
quieter where the instrument is doing something new. It will not ask the
reader to accept a physical identification before the build carries the
corresponding report.

It also retires a false obligation. The later episode chain is no longer
the order in which this book must be read. The first fifteen episodes
remain sacred as the frozen bootstrap. Episodes after that become a
quarry: some pieces may survive as sensors or reports; some may move to
the attic; some may remain useful only as provenance. The book follows
the instrument, not the sedimentary history of our attempts to
understand it.

The order now is:

\[
\begin{array}{c}
\hbox{sacred floor}
  \longrightarrow \hbox{sensor classes}
  \longrightarrow .999\hbox{ needle}\\[3pt]
  \longrightarrow \hbox{polar/Lorentz glue}
  \longrightarrow \hbox{verifier}.
\end{array}
\]

That is enough for a first count. It is enough for a frame. It is enough
to begin.

\section*{Vetting Ledger}

This preface claims no new theorem. It names the reading protocol for
the overhaul.

The floor blank, the finite needle, and the interval

\[
137.035999 \leq \alpha^{-1} \leq 137.036
\]

come from the current master specification. The local polar frame and
the alpha-particle, beta-particle, and gamma-particle channel glue come
from the current request-for-comment plan. The preface claims no global
polar Lorentz transform, no continuum limit, and no sharper numerical
value. The later episode files are not deleted by this preface; they are
marked for inventory and possible attic quarantine under the trim rule.

The next obligation is a build report that carries the three counted
dimensions, the channel closes, the ledger balance, and the finite box.
Until that report exists, the prose remains a gauge for the work, not a
substitute for it.
